% TOP_PROBLEM: {"id": 11100005, "title": "Major problems 5 — The chromatic splitting conjecture, which is considerably more complicated to state.", "category_id": 7, "category": {"id": 7, "name": "topology", "display_name": "Topology", "description": "Properties preserved under continuous deformations.", "slug": "topology", "order_index": 7, "created_at": "2026-07-31T15:26:25.671Z"}, "status": "partially_solved"}
% FIRST_POSTED: null
% ENTRY_KIND: statement_only
% !TeX program = lualatex
\documentclass[11pt]{article}
\usepackage[margin=25mm]{geometry}
\usepackage{fontspec}
\setmainfont{Latin Modern Roman}
\usepackage{amsmath,amssymb,mathtools}
\usepackage{unicode-math}
\setmathfont{Latin Modern Math}
\usepackage{xurl,hyperref}
\hypersetup{colorlinks=true,urlcolor=blue}
\setcounter{secnumdepth}{0}
\setlength{\parindent}{0pt}
\setlength{\parskip}{5pt}
\setlength{\emergencystretch}{3em}
\newcommand{\R}{\mathbb R}
\newcommand{\N}{\mathbb N}
\newcommand{\Z}{\mathbb Z}
\newcommand{\Q}{\mathbb Q}
\newcommand{\C}{\mathbb C}
\newcommand{\D}{\mathbb D}
\newcommand{\F}{\mathbb F}
\newcommand{\sref}[2]{\hyperlink{source:#1:#2}{[S#2]}}
\newcommand{\cataloguescope}[1]{#1}
\begin{document}
% UnsolvedMath ID 11100005; source catalogue code AMR-110-0005
\setcounter{section}{11100004}
\section{Major problems 5 — The chromatic splitting conjecture, which is considerably more complicated to state.}\label{11100005}
{\small UnsolvedMath ID 11100005 · Topology}
\subsection{Definitions and mathematical statement}

\paragraph{Shared notation.}$\mathbb N$, $\mathbb Z$, $\mathbb Q$, $\mathbb R$, and $\mathbb C$ denote the natural numbers, integers, rationals, reals, and complex numbers, respectively. All additional parameters and hypotheses are specified in the question.


\paragraph{Statement recorded in the supplied source.}
For every chromatic height $n$ and prime $p$, does the chromatic splitting conjecture hold for the $p$-local sphere? Write $L_n$ for localization at $E(n)$ and $L_{K(n)}$ for localization at Morava $K(n)$. The natural map $L_{n-1}S\to L_{n-1}L_{K(n)}S$ is predicted to exhibit $L_{n-1}S$ as a wedge summand, with the complementary spectrum given by the prescribed wedge of lower chromatic localizations and suspensions. This asks for the full splitting prediction, including its prescribed complementary summands.
\par\smallskip\textit{Formulation references:} \sref{11100005}{1}, \sref{11100005}{2}.
\subsection{Short English statement}
For every chromatic height $n$ and prime $p$, does the chromatic splitting conjecture hold for the $p$-local sphere? Write $L_n$ for localization at $E(n)$ and $L_{K(n)}$ for localization at Morava $K(n)$. The natural map $L_{n-1}S\to L_{n-1}L_{K(n)}S$ is predicted to exhibit $L_{n-1}S$ as a wedge summand, with the complementary spectrum given by the prescribed wedge of lower chromatic localizations and suspensions. This asks for the full splitting prediction, including its prescribed complementary summands.
\subsection{Sources}
\begin{itemize}
\item[\hypertarget{source:11100005:1}{S1}] ULAM.AI and the UnsolvedMath Contributors, UnsolvedMath: A Curated Collection of Open Mathematics Problems, record 11100005 (AMR-110-0005). Catalogue attribution under CC BY 4.0.
\url{https://huggingface.co/datasets/ulamai/UnsolvedMath}.
\item[\hypertarget{source:11100005:2}{S2}] Original problem formulation and mathematical context.
\url{https://www-users.cse.umn.edu/~tlawson/hovey/}.
\end{itemize}
\label{end:11100005}
\end{document}
